\begin{otherlanguage}{english}
\chapter*{\centering Abstract}
Retail networks must choose replenishment policies for heterogeneous demand series. This dissertation investigates whether demand characteristics can improve the cost--service trade-off relative to a global policy. The proposed Adaptive Inventory Policy Engine (AIPE) integrates data preparation, demand characterization, and simulation-based policy comparison. The research includes a pilot involving 15 stores in Paraíba, a historical study of 145 store--product series in Bahia, and a historical comparison using 2,408 M5 food series. Historical differences between policies depend on the cost criterion, service requirement, and experimental version. A new Bahia experiment evaluates the same 145-series cohort at two temporal origins, using seven policies including an additional control, three training seeds for stochastic methods, and separate validation and test blocks. All 4,350 fits are complete. Profile-based selection increases mean cost by 1.93\% and mean fill rate by 3.89 percentage points. Cases below the 0.70 target decrease from 119 to 100 among 283 series--origin cases with positive test demand. Fill rate is undefined in seven zero-demand cases, which are counted separately. The global choice changes between origins; at the second origin, both selectors recommend the same policy. The evidence does not establish general cost savings from profile-based selection. It identifies the cost of additional service and the temporal sensitivity of recommendations. The contribution is a formulation and traceable evaluation of contextual policy selection in Brazilian retail. Retrospective cohort membership and short, partially overlapping test windows limit generalization. The supervised Policy Selection Engine (PSE) remains to be validated, while the corrected M5 execution is a separate ongoing stage.

\vspace{2em}
\textbf{Keywords}: Inventory management; Intermittent demand; Policy selection; Service level; Simulation.
\end{otherlanguage}
